%% sections/methods.tex — Model and Methods  (revised after Round-1 polishing)
%% Word count target: ~850 words

\subsection{Transverse-field Ising model}
\label{sec:TFIM}

We study quantum quenches in the one-dimensional transverse-field Ising model (TFIM) on a
periodic chain of $N$ sites.
After a standard Jordan-Wigner transformation~\cite{Lieb1961}, the Hamiltonian takes the
free-fermion form
\begin{equation}
  H(g) = -\frac{J}{2}\sum_{j=1}^{N}\!\bigl(c_j^\dagger c_{j+1}^{\phantom{\dagger}}
  + c_j^\dagger c_{j+1}^\dagger + \mathrm{h.c.}\bigr)
  + g\sum_{j=1}^{N}(2n_j - 1),
  \label{eq:TFIM}
\end{equation}
where $c_j^\dagger$, $c_j$ are spinless fermionic operators satisfying canonical
anticommutation relations, $n_j = c_j^\dagger c_j$ is the site occupation, and
periodic boundary conditions ($c_{N+1} \equiv c_1$) are imposed.
The model undergoes an Ising-class quantum phase transition at the critical field
$\gc = J/2$~\cite{Sachdev2011}: for $g < \gc$ the ground state breaks the discrete Z2
symmetry (ferromagnetic phase), whereas for $g > \gc$ the symmetry is restored (paramagnetic
phase). Throughout this work we set $J = 1$, giving $\gc = 0.5$.

A remark on boundary conditions is in order.
On the periodic ring, the Jordan-Wigner transformation produces parity-dependent boundary
conditions: the even-parity (odd-parity) sector $H_\pm$ has
momenta $k = 2\pi n/N$ ($k = 2\pi(n+\tfrac{1}{2})/N$), $n = 0,\ldots, N-1$.
Because the initial state studied here has support in \emph{both} parity sectors
(see below), the time-evolved state is a coherent superposition of states governed by
these two spectrally distinct quadratic theories.
This sector mixing prevents reduction to a single-sector correlation-matrix
calculation~\cite{Fagotti2008}, and motivates the use of full exact diagonalization.

To probe the universality of our findings, we extend the analysis to the
anisotropic XY chain
\begin{equation}
  H_{\rm XY}(g,\gamma) = -\frac{J}{2}\sum_j\!\bigl(c_j^\dagger c_{j+1}^{\phantom{\dagger}}
  + \mathrm{h.c.}\bigr)
  - \frac{J(1-\gamma)}{2}\sum_j\!\bigl(c_j^\dagger c_{j+1}^\dagger + \mathrm{h.c.}\bigr)
  + g\sum_j(2n_j - 1),
  \label{eq:XY}
\end{equation}
where $\gamma \in [0,1]$ tunes the ratio of pairing to hopping: $\gamma = 0$ recovers the
TFIM~\eqref{eq:TFIM} while $\gamma = 1$ reduces to the XX model (vanishing pairing).
The critical field follows~\cite{Franchini2017}
\begin{equation}
  \gc(\gamma) = \frac{J}{2}\sqrt{1-\gamma},
  \label{eq:gc_XY}
\end{equation}
interpolating from $\gc = 0.5$ (TFIM) to $\gc = 0$ (XX model), consistently with the
known phase diagram of the anisotropic XY chain~\cite{Katsura1962}.

\subsection{Initial state and quantum quench}
\label{sec:quench}

We prepare the system in the polarized product state
\begin{equation}
  |\psi_\theta\rangle = \bigotimes_{j=1}^{N}\!\bigl[\cos\!\tfrac{\theta}{2}\,|0\rangle_j
  + \sin\!\tfrac{\theta}{2}\,|1\rangle_j\bigr], \quad \theta \in (0,\pi),
  \label{eq:psi0}
\end{equation}
where $|0\rangle_j$ and $|1\rangle_j$ denote empty and occupied fermionic sites,
and $\theta$ controls the initial polarization along the $z$-axis.
The system is then quenched: it evolves unitarily under the post-quench Hamiltonian
$H(\gf)$ for $t > 0$, starting from $|\psi_\theta\rangle$.
Product spin states of this form are the natural initial conditions in cold-atom quantum
simulation platforms~\cite{Bloch2008} and are particularly relevant because they are
\emph{not} eigenstates of the conserved Z2 symmetry.

The state~\eqref{eq:psi0} is a superposition of even and odd fermion-parity sectors:
$\langle(-1)^{N_f}\rangle_\theta = \cos^N\theta$, which vanishes as $N\to\infty$ for
$\theta \in (0,\pi/2)$ but is nonzero at any finite $N$.
As a consequence, computing the entanglement asymmetry (see below) requires the
off-diagonal elements of the reduced density matrix across parity sectors, which
are not encoded in two-point correlation functions alone~\cite{Fagotti2008}.
This prevents the use of Gaussian correlation-matrix methods~\cite{Ares2023} and
necessitates exact diagonalization at finite system size.

\subsection{Z2 symmetry and entanglement asymmetry}
\label{sec:EA}

The Hamiltonian~\eqref{eq:TFIM} commutes with the fermion parity operator
\begin{equation}
  Q = (-1)^{\hat{N}_f} = \prod_{j=1}^{N}(1-2n_j),
  \label{eq:Q}
\end{equation}
which in the original spin language corresponds to $\prod_j\sigma^z_j$ (the global
magnetization parity), and is distinct from the Ising spin-flip $\prod_j\sigma^x_j$.
Both ground-state phases have definite fermion parity, so neither breaks $Q$; the initial
product state~\eqref{eq:psi0} does, via $\langle Q\rangle_\theta = \cos^N\theta$.
Note that $Q$ is \emph{not} a U(1) particle-number symmetry: the pairing terms
$c_j^\dagger c_{j+1}^\dagger$ in~\eqref{eq:TFIM} break U(1) explicitly, leaving only
the discrete Z2 generated by $Q$.

To quantify how much of this initial symmetry breaking is retained in a
subsystem $A = \{1,\ldots,\NA\}$, we use the \emph{entanglement asymmetry}~\cite{Ares2023}
\begin{equation}
  \DS_A(t) = S(\tilde\rho_A(t)) - S(\rho_A(t)) \geq 0,
  \label{eq:EA}
\end{equation}
where $\rho_A = \mathrm{Tr}_B\,|\psi(t)\rangle\langle\psi(t)|$ is the reduced density matrix
of $A$, $\tilde\rho_A = (\rho_A + Q_A\rho_A Q_A)/2$ is its Z2-projected counterpart with
$Q_A = \prod_{j\in A}(1-2n_j)$, and $S(\rho) = -\mathrm{Tr}[\rho\ln\rho]$ is the von
Neumann entropy. The quantity $\DS_A(t)$ vanishes if and only if $[\rho_A, Q_A] = 0$,
signalling complete Z2 restoration in $A$, and is bounded by
$\DS_A \leq \ln 2$ for any state (since $|G| = 2$)~\cite{Ares2023}.

For the product state~\eqref{eq:psi0}, the reduced density matrix $\rho_A(0)$ is a pure
state ($S(\rho_A) = 0$), and the Z2-projected version
$\tilde\rho_A(0) = (\rho_A + Q_A\rho_A Q_A)/2$ has exactly two non-zero eigenvalues
$p_\pm = (1 \pm \langle Q_A\rangle_\theta)/2 = (1 \pm \cos^{\NA}\theta)/2$.
Inserting these into~\eqref{eq:EA} yields the exact initial formula
\begin{equation}
  \DS_A(0) = \hbin\!\!\left(\frac{1+\cos^{\NA}\!\theta}{2}\right),
  \label{eq:DS0}
\end{equation}
where $\hbin(p) = -p\ln p - (1-p)\ln(1-p)$ is the binary entropy function.
This result saturates the universal bound $\DS_A \leq \ln 2$ at $\theta = \pi/2$ (maximal
symmetry breaking, $\langle Q_A\rangle = 0$) and vanishes at $\theta = 0, \pi$
(Z2-symmetric initial states).
We verified Eq.~\eqref{eq:DS0} against exact numerical diagonalization to
relative errors below $10^{-14}$ for all $\theta \in (0,\pi)$ and $\NA = 1,\ldots,8$.

\subsection{Quantum Mpemba effect and crossing time}
\label{sec:QME}

We say a pair $(\theta_1, \theta_2)$ with $\DS_A(0;\theta_1) > \DS_A(0;\theta_2)$ exhibits
the \emph{quantum Mpemba effect} (QME) if there exists a finite time $\tM > 0$ such that
$\DS_A(\tM;\theta_1) = \DS_A(\tM;\theta_2)$, with the more-symmetry-broken state
subsequently falling below the less-broken one.
In direct analogy with the classical Mpemba effect — in which a hotter system can
equilibrate to a cooler environment faster than a less hot one~\cite{Mpemba1969} — the
\emph{more broken} initial state restores symmetry \emph{faster}, crossing the less
broken state in EA at $\tM$.
We define $\tM$ as the earliest such crossing and detect it by linear interpolation
between time steps. Pairs with initial EA gap $|\DS_A(\theta_1,0) - \DS_A(\theta_2,0)|
< 0.05$ are excluded as insufficiently distinguished.

\subsection{Dynamical quantum phase transitions}
\label{sec:DQPT}

Dynamical quantum phase transitions (DQPTs)~\cite{Heyl2013,Heyl2018} manifest as
non-analytic cusps in the Loschmidt return rate
\begin{equation}
  f(t) = -\frac{1}{N}\ln\bigl|\mathcal{L}(t)\bigr|^2, \qquad
  \mathcal{L}(t) = \langle\psi_\theta|e^{-iH(\gf)t}|\psi_\theta\rangle.
  \label{eq:losch}
\end{equation}
For quenches into the paramagnetic phase ($\gf > \gc$), the return rate develops cusps at
times $t_n^* = (2n+1)T^*/2$, $n = 0, 1, 2, \ldots$, where the DQPT period
is~\cite{Heyl2013}
\begin{equation}
  T^* = \frac{\pi}{\gf - \gc}.
  \label{eq:Tstar}
\end{equation}
In finite systems, we identify DQPTs as prominent local maxima of $f(t)$ exceeding a
threshold $f > 0.05$.
To characterize the relationship between QME crossings and the DQPT structure, we
introduce the \emph{fractional crossing time}
\begin{equation}
  \phi = \frac{t_M \bmod T^*}{T^*} \in [0,1),
  \label{eq:phi}
\end{equation}
which measures where within a single DQPT period the crossing falls.
DQPT singularities correspond to $\phi = 0.5$; a distribution concentrated near $\phi = 0.5$
would indicate QME crossings locked to DQPT times.

\subsection{Numerical methods}
\label{sec:numerics}

We compute the time-evolved wavefunction $|\psi(t)\rangle = e^{-iH(\gf)t}|\psi_\theta\rangle$
in the full $2^N$-dimensional Hilbert space using the Krylov-subspace
algorithm~\cite{AlMohy2011} (\textsc{SciPy}~\cite{Virtanen2020} routine
\texttt{expm\_multiply}), which avoids explicit construction of the time-evolution operator
and scales as $\mathcal{O}(2^N \cdot N_{\rm Krylov})$ per timestep.
The entanglement asymmetry~\eqref{eq:EA} is evaluated by explicit partial trace to form
the $2^{\NA}\times 2^{\NA}$ reduced density matrix $\rho_A$, followed by Z2 dephasing and
numerical diagonalization.

System sizes range from $N = 8$ to $N = 18$ (Hilbert-space dimension $2^{18} \approx
2.6\times 10^5$) with subsystem size $\NA = 4$ for the main analysis; the sensitivity to
$\NA$ is studied separately for $\NA \in \{2, 3, 4, 5, 6\}$ at $N = 18$
(Sec.~\ref{sec:xy}). Time evolution uses $n_t = 150$--$200$ uniformly spaced points over
$t_{\max} = 15$--$20\, J^{-1}$, with timestep $\delta t \lesssim 0.1 J^{-1}$.
Finite-size effects on $\tM$ are analysed in Sec.~\ref{sec:fss}, where we demonstrate
via linear extrapolation in $1/N$ that $\tM$ converges smoothly to a finite
thermodynamic-limit value for all parameter pairs studied, confirming that $N = 18$ sits
in the scaling regime for our observables.
